\documentclass[11pt,a4paper]{article}

\usepackage[utf8]{inputenc} % allow utf-8 input
\usepackage[T1]{fontenc}    % use 8-bit T1 fonts
\usepackage{hyperref}       % hyperlinks
\usepackage{url}            % simple URL typesetting
\usepackage{booktabs}       % professional-quality tables
\usepackage{amsfonts}       % blackboard math symbols
\usepackage{nicefrac}       % compact symbols for 1/2, etc.
\usepackage{microtype}      % microtypography
\usepackage{xcolor}         % colors
\usepackage{graphicx}       % graphics
\usepackage[normalem]{ulem} % strikethrough
\title{Global transcriptional response of}

\begin{document}

\maketitle

Yiquan Hu, Yarong Miao

\section{Abstract}

\subsection{Background}

Environmental modulation of gene expression in Yersinia pestis is critical for its life style and pathogenesis. Using cDNA microarray technology, we have analyzed the global gene expression of this deadly pathogen when grown under different stress conditions in vitro.

\subsection{Results}

To provide us with a comprehensive view of environmental modulation of global gene expression in Y. pestis, we have analyzed the gene expression profiles of 25 different stress conditions. Almost all known virulence genes of Y. pestis were differentially regulated under multiple environmental perturbations. Clustering enabled us to functionally classify co-expressed genes, including some uncharacterized genes. Collections of operons were predicted from the microarray data, and some of these were confirmed by reverse-transcription polymerase chain reaction (RT-PCR). Several regulatory DNA motifs, probably recognized by the regulatory protein Fur, PurR, or Fnr, were predicted from the clustered genes, and a Fur binding site in the corresponding promoter regions was verified by electrophoretic mobility shift assay (EMSA).

\subsection{Conclusion}

The comparative transcriptomics analysis we present here not only benefits our understanding of the molecular determinants of pathogenesis and cellular regulatory circuits in Y. pestis, it also serves as a basis for integrating increasing volumes of microarray data using existing methods.

\section{Introduction}

Yersinia pestis is the causative agent of plague, one of the most devastating           infectious diseases in human history.

\end{document}